\documentclass[11pt]{article}

\usepackage[english]{babel}
\usepackage[letterpaper,top=2cm,bottom=2cm,left=2.7cm,right=2.7cm,marginparwidth=1.6cm]{geometry}
\usepackage{amsmath,amssymb}
\usepackage{graphicx}
\usepackage{booktabs}
\usepackage[colorlinks=true, allcolors=blue]{hyperref}

\newcommand{\g}[1]{g^{(#1)}}
\newcommand{\vx}{\mathbf{x}}
\newcommand{\vy}{\mathbf{y}}
\newcommand{\rd}{\,\mathrm{d}}
\newcommand{\kB}{k_{\mathrm B}}
\newcommand{\Dt}{D_{\mathrm t}}
\newcommand{\Dr}{D_{\mathrm r}}

\title{Extended fitting methods for multi-angle, polarized photon correlation
spectroscopy:\\ generalizing and relaxing the assumptions of CORENN}
\author{Michael L. Whittaker}
\date{}

\begin{document}
\maketitle

\begin{abstract}
The CORENN algorithm recovers a particle-size distribution (PSD) from a single
dynamic-light-scattering (DLS) correlation curve by a non-linear,
noise-weighted, Tikhonov-regularized, non-negativity-constrained inversion of
the Siegert model. It is, however, built on four assumptions that are violated
by the photon-correlation-tomography (PCT) measurements we wish to interpret:
(i)~a single scattering angle, so that the intensity weighting of each size
class is treated as $q$-independent; (ii)~optically isotropic, spherical
scatterers undergoing purely \emph{translational} diffusion; (iii)~a correlation
curve supplied entirely by the hardware correlator; and (iv)~a parametric noise
model in place of measured replicate statistics. This document develops a
unified inverse-problem framework that relaxes each assumption. We (1)~promote
the unknown to a $q$-consistent distribution and fold the angular dependence of
the scattering through the particle form factor $P(qR)$; (2)~add the
\emph{rotational} degree of freedom probed by depolarized (VH) scattering,
yielding a joint distribution over translational and rotational dynamics
$f(\Dt,\Dr)$ from which size \emph{and} shape/anisotropy are recovered;
(3)~extend the usable lag range by a graceful, uncertainty-weighted merge of the
hardware correlogram with a software correlation of the photon count-rate trace;
and (4)~replace the noise model with a replicate-calibrated, Schätzel-shaped
weighting. The resulting estimator is a single regularized non-linear weighted
least-squares problem over a non-negative distribution, of which standard CORENN
is the spherical, single-angle, translational-only special case.
\end{abstract}

\section{Background: the CORENN forward model}
The measured intensity correlation function (ICF) is related to the field
correlation function (FCF) by the Siegert relation
\begin{equation}\label{eq:siegert}
  \g2(q,\tau) - 1 = \beta\,\bigl|\g1(q,\tau)\bigr|^2 ,
\end{equation}
with coherence factor $\beta\le 1$. For a dilute, ergodic suspension of
spherical scatterers the FCF is the Laplace transform of the intensity-weighted
distribution of decay rates,
\begin{equation}\label{eq:g1_corenn}
  \g1(q,\tau) = \int_0^\infty I(R)\,e^{-q^2 D(R)\,\tau}\rd R,
  \qquad
  D(R) = \frac{\kB T}{6\pi\eta R},
  \qquad \int_0^\infty I(R)\rd R = 1 ,
\end{equation}
where $I(R)$ is the fraction of scattered intensity arising from particles of
hydrodynamic radius $R$, $D(R)$ is the Stokes--Einstein translational diffusion
coefficient, $\eta$ the solvent viscosity and $q = (4\pi n/\lambda_0)\sin(\theta/2)$
the scattering-vector magnitude. CORENN discretizes $I(R)$ on a cubic B-spline
basis $\{\Phi_j\}$,
\begin{equation}\label{eq:bspline}
  I(R) \approx \sum_{j=1}^{m} x_j\,\Phi_j(R),\qquad x_j\ge 0,
\end{equation}
so that, defining the kernel
$K_{ij}=\int \Phi_j(R)\,e^{-q^2 D(R)\tau_i}\rd R$, the model reads
$\g2(\tau_i)-1 = \beta\,(K\vx)_i^2$. The PSD is then the solution of the
regularized non-linear weighted least-squares (NLWLS) problem
\begin{equation}\label{eq:corenn}
  \hat\vx(\lambda) = \arg\min_{\vx\ge 0}\;
   \sum_i w_i^2\bigl[\,\hat g_i - \beta\,(K\vx)_i^2 - b\,\bigr]^2
   \;+\; \lambda\,\|L\vx\|^2 ,
\end{equation}
with data weights $w_i = 1/\sigma_i$, a free baseline $b$, a curvature penalty
$L$ (the second-derivative operator $\int [I''(R)]^2\rd R = \vx^\top L^\top L\,\vx$),
and the regularization weight $\lambda$ chosen by generalized cross validation
(GCV). The defining feature relative to CONTIN is that the \emph{square} in
Eq.~\eqref{eq:corenn} is retained: one does not linearize by taking
$\sqrt{(\hat g-1)/\beta}=K\vx$, which would distort the (heteroscedastic) noise.

\paragraph{Assumptions to be relaxed.}
Equations~\eqref{eq:siegert}--\eqref{eq:bspline} embed: a single $q$ with
$q$-independent weighting $I(R)$ (\S\ref{sec:multiangle}); spherical, optically
isotropic particles with only translational diffusion (\S\ref{sec:polarized});
a hardware-supplied ICF over a fixed lag window (\S\ref{sec:hybrid}); and a
parametric noise model for $w_i$ (\S\ref{sec:noise}).

\section{Multi-angle inversion with explicit form-factor weighting}
\label{sec:multiangle}
The intensity scattered by a particle of radius $R$ at scattering vector $q$ is
not $q$-independent: for a homogeneous sphere in the Rayleigh--Gans--Debye
limit it scales as $R^6 P(qR)$, with the form factor
\begin{equation}\label{eq:formfactor}
  P(qR) = \left[\,\frac{3\bigl(\sin qR - qR\cos qR\bigr)}{(qR)^3}\,\right]^2 ,
  \qquad P(qR)\to 1 \text{ as } qR\to 0 .
\end{equation}
Writing the number distribution as $N(R)$, the intensity weighting at angle $q$
is $I(R,q) \propto N(R)\,R^6\,P(qR)$. The crucial observation is that the
$R^6$ factor is $q$-independent, so we may define the \emph{angle-extrapolated
intensity distribution}
\begin{equation}
  I_0(R) \;\equiv\; \frac{N(R)\,R^6}{\int N(R')\,R'^6\rd R'},
\end{equation}
i.e.\ the intensity PSD one would measure as $q\to 0$. The per-angle FCF is then
\begin{equation}\label{eq:g1_multiangle}
  \g1(q,\tau) = \frac{\displaystyle\int I_0(R)\,P(qR)\,e^{-q^2 D(R)\,\tau}\rd R}
                     {\displaystyle\int I_0(R)\,P(qR)\rd R} ,
\end{equation}
where the denominator restores $\g1(q,0)=1$. This relaxes CORENN's implicit
assumption that the same $I(R)$ applies at every angle: the only $q$-dependent
re-weighting of a fixed $I_0(R)$ is the bounded, $O(1)$ factor $P(qR)$, which is
numerically benign (unlike the $\sim$20-decade dynamic range of $R^6$ itself).

Discretizing $I_0(R)=\sum_j x_j\Phi_j(R)$ gives, for each angle $q^{(a)}$, a
kernel
$K^{(a)}_{ij}=\int \Phi_j(R)\,P(q^{(a)}R)\,e^{-(q^{(a)})^2 D(R)\tau_i}\rd R$ and a
model $\g2_a(\tau_i)-1 = A_a\,(K^{(a)}\vx)_i^2$, with a per-angle amplitude
$A_a$ (coherence factor divided by the squared normalization) that can be
profiled out in closed form,
\begin{equation}
  A_a(\vx) =
  \frac{\sum_i w_{a,i}^2\,\hat g_{a,i}\,(K^{(a)}\vx)_i^2}
       {\sum_i w_{a,i}^2\,(K^{(a)}\vx)_i^4}.
\end{equation}
The joint estimator solves a single non-negative, regularized NLWLS over the
$q$-independent coefficients $\vx$ using \emph{all} angles simultaneously,
\begin{equation}\label{eq:joint_multiangle}
  \hat\vx(\lambda) = \arg\min_{\vx\ge 0}\;
  \sum_{a}\sum_i w_{a,i}^2\bigl[\hat g_{a,i} - A_a(\vx)\,(K^{(a)}\vx)_i^2\bigr]^2
  + \lambda\,\|L\vx\|^2 .
\end{equation}
Because the form factor develops minima and lobes once $qR\gtrsim\pi$, the
angular data carry information about the \emph{large}-particle tail that a
single back-scattering angle cannot resolve; conversely the small-particle
mode, whose intensity weight is suppressed by $R^6$, remains weakly constrained
(an intrinsic DLS limitation, not a defect of the inversion). The number and
volume distributions follow from $N(R)\propto I_0(R)/R^6$ and
$V(R)\propto I_0(R)/R^3$.

\section{Polarized DLS: adding rotational diffusion}
\label{sec:polarized}
For optically anisotropic or non-spherical particles the scattered field has a
depolarized component, and the dynamics carry information about rotation.
Decompose the scattering into polarized (VV; incident and detected
polarizations both vertical) and depolarized (VH; crossed) channels. To leading
order in the optical anisotropy, the field correlation functions are
\begin{align}
  \g1_{\mathrm{VV}}(q,\tau)
    &= \int f(\Dt)\,e^{-q^2 \Dt\,\tau}\rd\Dt
       \;+\; \tfrac{4}{3}\,\delta^2\!\!\int f(\Dt,\Dr)\,
       e^{-(q^2\Dt + 6\Dr)\tau}\rd\Dt\rd\Dr ,
       \label{eq:g1VV}\\[2pt]
  \g1_{\mathrm{VH}}(q,\tau)
    &= \int f(\Dt,\Dr)\,e^{-(q^2\Dt + 6\Dr)\,\tau}\rd\Dt\rd\Dr ,
       \label{eq:g1VH}
\end{align}
where $\Dt$ and $\Dr$ are the translational and rotational diffusion
coefficients, $\delta$ is the (small) optical-anisotropy factor, and
$f(\Dt,\Dr)$ is the joint distribution of dynamics weighted by scattered
intensity. The depolarized channel \eqref{eq:g1VH} decays with the
\emph{enhanced} rate $\Gamma_{\mathrm{VH}} = q^2\Dt + 6\Dr$: at low $q$ it is
dominated by the $q$-independent rotational term $6\Dr$, providing direct access
to rotation. For spheres $\delta = 0$ and the VH channel vanishes, recovering
the CORENN picture; the polarized framework is therefore a strict generalization.

\paragraph{From dynamics to geometry.}
The two diffusion coefficients constrain two geometric parameters. For a
spherical equivalent, $\Dt = \kB T/6\pi\eta R_{\mathrm t}$ and
$\Dr = \kB T/8\pi\eta R_{\mathrm r}^3$ define translational and rotational
hydrodynamic radii; their ratio
\begin{equation}
  \rho \;\equiv\; \frac{\Dr}{\Dt}\,R_{\mathrm t}^2
  \;=\; \frac{3}{4}\left(\frac{R_{\mathrm t}}{R_{\mathrm r}}\right)^{\!3}
\end{equation}
is a dimensionless shape descriptor (unity for a sphere, $>1$ for elongated
particles). More generally, for an axisymmetric ellipsoid of semi-axes $(a,b)$
the Perrin friction factors map $(\Dt,\Dr)$ to the aspect ratio $p=a/b$ and
volume, so the recovered $f(\Dt,\Dr)$ becomes a joint distribution over size and
shape.

\paragraph{Joint VV/VH multi-angle inverse problem.}
Combining \S\ref{sec:multiangle} and the present section, the unknown is a
non-negative joint distribution $f(R,\,s)$ over hydrodynamic size $R$ and a
shape coordinate $s$ (e.g.\ aspect ratio), discretized on a tensor B-spline
basis $f(R,s)=\sum_{jk} x_{jk}\,\Phi_j(R)\,\Psi_k(s)$, $x_{jk}\ge 0$. Each
(angle, polarization) channel $c=(q^{(a)},\mathrm{pol})$ contributes a kernel
\begin{equation}\label{eq:joint_kernel}
  K^{(c)}_{i,jk} = \int\!\!\int W_c(R,s,q^{(a)})\,\Phi_j(R)\,\Psi_k(s)\,
  \exp\!\bigl[-\Gamma_c(R,s,q^{(a)})\,\tau_i\bigr]\rd R\rd s ,
\end{equation}
with channel-specific weight $W_c$ (intensity $\times$ form factor $\times$
polarization selection rule) and decay rate
\begin{equation}
  \Gamma_c =
  \begin{cases}
    q^2\,\Dt(R) & \text{VV, leading order},\\[2pt]
    q^2\,\Dt(R) + 6\,\Dr(R,s) & \text{VH}.
  \end{cases}
\end{equation}
All channels are inverted jointly,
\begin{equation}\label{eq:joint_full}
  \hat\vx(\lambda) = \arg\min_{\vx\ge 0}
  \sum_{c}\sum_i w_{c,i}^2\bigl[\hat g_{c,i} - A_c(\vx)\,(K^{(c)}\vx)_i^2\bigr]^2
  + \lambda\,\vx^\top \Omega\,\vx ,
\end{equation}
where $\Omega$ is a separable curvature penalty enforcing smoothness in both $R$
and $s$. Equation~\eqref{eq:joint_full} reduces to standard CORENN
\eqref{eq:corenn} for a single VV angle and $s$-independent (spherical)
kinematics, and to the multi-angle problem \eqref{eq:joint_multiangle} when the
VH channels and shape coordinate are dropped.

\section{Hybrid hardware/software correlation for an extended lag range}
\label{sec:hybrid}
Resolving both the fast rotational dynamics of \S\ref{sec:polarized} and the
slow dynamics of large aggregates requires a correlation function spanning many
more decades in $\tau$ than any single estimator provides. The hardware
correlator is accurate from $\sim$ns to the end of the main relaxation but, at
long lags, suffers triangular-averaging distortion, few independent
decorrelation windows, and irreproducible slow ``artifacts.'' Meanwhile the
instrument streams a photon count-rate trace $n(t)$ binned every few
milliseconds for the whole acquisition. Software-correlating that trace --- and
ensemble-averaging across replicates --- yields a superior estimate of the slow
dynamics. We use the symmetric-normalization, cross-correlation estimator
\begin{equation}\label{eq:soft_corr}
  \g2_{\mathrm{sw}}(\tau_k) =
  \frac{\bigl\langle n_A(t_i)\,n_B(t_i+\tau_k)\bigr\rangle_i}
       {\bigl\langle n_A(t_i)\bigr\rangle_i\,\bigl\langle n_B(t_i+\tau_k)\bigr\rangle_i},
\end{equation}
with the direct and delayed monitor averages taken over the same overlapping
segment (suppressing drift bias) and channels $A\times B$ cross-correlated to
remove detector afterpulsing and dead-time.

\paragraph{Graceful, uncertainty-weighted merge.}
Rather than hard-truncating the hardware curve at an arbitrary lag, the two
estimators are blended by inverse-variance weighting on a common log-lag grid:
\begin{equation}\label{eq:merge}
  \g2_{\mathrm{merge}}(\tau) - 1 =
  \frac{\sigma_{\mathrm{hw}}^{-2}\,[\g2_{\mathrm{hw}}-1]
        + \sigma_{\mathrm{sw}}^{-2}\,[\g2_{\mathrm{sw}}-1]}
       {\sigma_{\mathrm{hw}}^{-2}+\sigma_{\mathrm{sw}}^{-2}},
  \qquad
  \sigma_{\mathrm{merge}}^{-2}=\sigma_{\mathrm{hw}}^{-2}+\sigma_{\mathrm{sw}}^{-2}.
\end{equation}
The software amplitude is first matched to the hardware over a narrow window
straddling the handoff (both measure $\beta\,|\g1|^2$ for the same field, but
with different coherence factors). The expected uncertainties encode where each
estimator is trustworthy:
\begin{align}
  \sigma_{\mathrm{hw}}(\tau)
    &= \sigma^{\mathrm{base}}_{\mathrm{hw}}(\tau)\,
       \Bigl[1 + (\tau/\tau_{\mathrm{distort}})^{p}\Bigr], \\
  \sigma_{\mathrm{sw}}(\tau)
    &= \frac{\sigma^{\mathrm{base}}_{\mathrm{sw}}(\tau)}{\sqrt{N_{\mathrm{rep}}}}\,
       \Bigl[1 + (\,m\,\Delta t_{\mathrm{bin}}/\tau)^{p}\Bigr],
\end{align}
so the hardware is down-weighted only beyond the distortion lag
$\tau_{\mathrm{distort}}$, while the software is down-weighted within a few bin
times $m\,\Delta t_{\mathrm{bin}}$ of its shortest lag. The crossover then falls
naturally just after the main decay, and $\sigma_{\mathrm{merge}}(\tau)$ is
exactly the weight $w_i=1/\sigma_i$ required by the inversions of
\S\S\ref{sec:multiangle}--\ref{sec:polarized}.

\section{Noise model: replicate-calibrated Schätzel weighting}
\label{sec:noise}
The statistical noise on the ICF estimator grows at long lag because the number
of independent decorrelation windows in a run of duration $T$ falls as
$M(\tau)\sim T/\tau$ (Schätzel 1990; Biganzoli \& Ferri 2018). A robust,
numerically stable model for the variance of $\g2-1$ is
\begin{equation}\label{eq:schatzel}
  \sigma^2(\tau) = \sigma_{\mathrm{floor}}^2
  + g_{\mathrm a}\,\frac{\bigl[1+\beta\,\g1(\tau)^2\bigr]^2}{M(\tau)},
  \qquad M(\tau)=\max\!\bigl(T/\tau,\,1\bigr),
\end{equation}
where $\sigma_{\mathrm{floor}}$ is the short-lag (shot/quantization) level and
$g_{\mathrm a}=O(1)$ a self-beating gain. When replicate curves are available
(the PCT data provide five per measurement), we take the \emph{shape} of
$\sigma(\tau)$ from Eq.~\eqref{eq:schatzel} but calibrate its absolute level to
the measured short-lag replicate scatter, and never under-weight a lag the
replicates flag as noisy: $\sigma \leftarrow \max(\sigma_{\text{model}},
\sigma_{\text{empirical}})$. This combines the long-lag growth that the
hardware artifacts demand with the empirical fidelity of direct replicate
statistics, and is what allows the inversion to ignore the long-lag noise floor
rather than fit it with a spurious slow population.

\section{Summary of the generalization}
Table~\ref{tab:generalization} situates the four extensions relative to standard
CORENN. Each is implemented in the accompanying \texttt{corenn.py} (inversion,
form-factor multi-angle, replicate/Schätzel weighting, uncertainty bands) and
\texttt{hybrid\_correlator.py} (software correlation and graceful merge); the
polarized joint inversion \eqref{eq:joint_full} is the natural next step within
the same NLWLS machinery.

\begin{table}[h]
\centering
\begin{tabular}{@{}lll@{}}
\toprule
Assumption in CORENN & Generalization & Mechanism \\
\midrule
Single angle, $q$-indep.\ $I(R)$ & Multi-angle $I_0(R)$ & form factor $P(qR)$, Eq.~\eqref{eq:joint_multiangle} \\
Spherical, isotropic, transl.\ only & Rotational dynamics $f(\Dt,\Dr)$ & VV/VH channels, Eqs.~\eqref{eq:g1VV}--\eqref{eq:joint_full} \\
Hardware ICF, fixed lag window & Hybrid hardware+software & inverse-variance merge, Eq.~\eqref{eq:merge} \\
Parametric noise model & Replicate-calibrated weights & Schätzel shape, Eq.~\eqref{eq:schatzel} \\
\bottomrule
\end{tabular}
\caption{Each row relaxes one CORENN assumption while preserving the
non-negative, regularized, non-linear weighted least-squares structure; CORENN
is the common special case.}
\label{tab:generalization}
\end{table}

\paragraph{Relation to existing PCT theory.}
This document extends the basic FCF/ICF and Siegert treatment of
\texttt{dls\_theory.tex}; Eqs.~\eqref{eq:siegert}--\eqref{eq:g1_corenn} here are
the polydisperse generalization of its Eq.~(13), and \S\ref{sec:polarized}
supplies the rotational content absent from the spherical, translational-only
model used there.

\end{document}
